\scanpage{80}
\mspage{73}
% archive sheet number: 79
b) $I$ catégorie
\diagram{
\begin{tikzcd}[ampersand replacement=\&,column sep=tiny]
0\&\unsure{-|\to}\&1\&\unsure{-|\to}\&2\&\unsure{-|\to}\&3\arrow[r]\&\ldots
\end{tikzcd}
}
i.e. la catégorie ordonnée associée à l'ens ordonné $\mathbb{N}$, pour la relation d'ordre ordinaire $m\leq n$. Donc $M(I)$ est formé des systèmes \uline{projectifs} dans $M$, indexés par $\mathbb{N}$. \deleted{\unsure{On}} On prouve facilement \inserted{(en procédant comme pour le cas du lemme p. 63)} que $M(I)_f$ est formé des syst. projectifs
\diagram{
\begin{tikzcd}[ampersand replacement=\&,column sep=small]
X_0\&X_1\arrow[l]\&X_2\arrow[l]\&\ldots\arrow[l]\&X_n\&X_{n+1}\arrow[l,"\unsure{p_n|\rho_n}"']
\end{tikzcd}
}
tels que les $X_i$ soient dans $M_f$, et les $p_n$ dans $F$. La question est donc: si $X_\bullet,Y_\bullet$ sont fibrants et $f_\bullet:X_\bullet\to Y_\bullet$ est tel que $f_i\in W_f$ pour $i\in\mathbb{N}$, alors a-t-on $\varprojlim f_i:X_\infty\to Y_\infty$ dans $W$?

Si \deleted{\unsure{la}} la réponse est affirmative \inserted{pour $\mathbb{N}$}, (pour un $W$ et une quillénisation \unsure{donné|donnés|donnée}), elle l'est aussi pour \deleted{toute} catégorie discrète dénombrable, comme on le voit \unsure{aussitôt}. Mais je ne vois pas comment déduire, $=$ \unsure{ça des axiomes|des axiomes|de ces axiomes}, disons pour un triplet de Quillen clos (closed model category) tel que $M$ soit de plus stable par $\varinjlim$ filtrantes, que celles-ci soient exactes, qu'il y ait
